\documentclass[a4paper,12pt]{article}
\usepackage[T2A]{fontenc}
\usepackage[utf8]{inputenc}
\usepackage[russian]{babel}
\usepackage{amsmath,amssymb,mathtools}
\usepackage{geometry}
\usepackage{microtype}
\usepackage{booktabs,tabularx,array}
\usepackage{enumitem}
\usepackage{needspace}
\usepackage{fancyhdr}
\usepackage{tikz}
\usepackage{xcolor}
\usepackage[colorlinks=true,linkcolor=tealshade,urlcolor=tealshade]{hyperref}
\geometry{left=24mm,right=24mm,top=23mm,bottom=25mm}
\definecolor{tealshade}{HTML}{007D81}
\definecolor{bodyink}{HTML}{303438}
\definecolor{softred}{HTML}{A33332}
\definecolor{softblue}{HTML}{235F8B}
\color{bodyink}
\setlength{\parindent}{0pt}
\setlength{\parskip}{7pt plus 2pt minus 1pt}
\setlength{\headheight}{17pt}
\linespread{1.13}
\setlist[itemize]{leftmargin=18pt,itemsep=3pt,topsep=3pt}
\pagestyle{fancy}
\fancyhf{}
\fancyhead[L]{\small Ревью домашней работы}
\fancyhead[R]{\small Ярослав Верещагин}
\renewcommand{\headrulewidth}{0.45pt}
\fancyfoot[C]{\small\thepage}
\renewcommand{\footrulewidth}{0pt}
\newcommand{\tealheading}[1]{\par\needspace{4\baselineskip}\vspace{9pt}\noindent{\large\color{tealshade}#1}\par\vspace{3pt}\noindent\textcolor{tealshade}{\rule{\linewidth}{0.35pt}}\par\vspace{2pt}}
\newcommand{\subfield}[1]{\par\needspace{3\baselineskip}\textbf{#1}\quad}
\newcommand{\errorfield}{\par\needspace{3\baselineskip}\textcolor{softred}{\textbf{Ошибка:}}\quad}
\newcommand{\ideafield}{\par\needspace{3\baselineskip}\textcolor{softblue}{\textbf{Ключевая идея:}}\quad}
\newcommand{\resultfield}{\par\needspace{3\baselineskip}\textcolor{tealshade}{\textbf{Итог:}}\quad}
\newcommand{\taskheading}[2]{\newpage\phantomsection\label{task:#1}\tealheading{Задание №\,#1. #2}}
\hypersetup{pdftitle={Ревью домашней работы — 16 ноября 2026},pdfauthor={Лёвин Артём Александрович}}
\begin{document}

% Титульная страница по академическому референсу
\thispagestyle{empty}
\begin{center}
\vspace*{30mm}
{\large\color{tealshade} МАТЕМАТИКА\par}
\vspace{19mm}
{\Huge Ревью домашней работы\par}
\vspace{12mm}
{\LARGE\color{tealshade} Геометрия и алгебра\par}
\vspace{5mm}
{\large Диагонали ромба и разложение на множители\par}
\vfill
{\large Ярослав Верещагин\par}
\vspace{10mm}
{\large 16 ноября 2026 г.\par}
\vfill
{Лёвин Артём Александрович\par}
{эксклюзивно для Ярослава Верещагина\par}
\vspace{15mm}
\begin{tikzpicture}[x=1cm,y=1cm]
  \draw[tealshade,line width=0.85pt] (0,0) .. controls (1.2,.46) and (2.15,-.42) .. (3.6,0)
    .. controls (4.9,.44) and (5.9,-.42) .. (7.25,0)
    .. controls (8.35,.33) and (9.15,-.17) .. (9.75,0);
  \fill[tealshade] (3.6,0) circle (1.5pt);
\end{tikzpicture}
\end{center}
\clearpage

\tealheading{Сводка ошибок}
В таблице приведены только те задания, в которых требуется исправление. Нажатие на номер переводит к подробному разбору.
\par\vspace{8pt}
{\small
\renewcommand{\arraystretch}{1.35}
\begin{tabularx}{\linewidth}{@{}>{\centering\arraybackslash}p{11mm} p{26mm} X p{23mm} p{25mm}@{}}
\toprule
\textbf{№} & \textbf{Тема} & \textbf{Суть ошибки} & \textbf{Ваш ответ} & \textbf{Правильный\newline ответ}\\
\midrule
\hyperref[task:9]{\textbf{9}} & Ромб & Пополам разделён острый угол вместо тупого. & $35^\circ$ & $55^\circ$\\
\hyperref[task:10]{\textbf{10}} & Разложение на множители & При раскрытии скобок изменён знак; не проверен второй множитель. & $4;\;-3$ & $-4;\;3$\\
\bottomrule
\end{tabularx}
}
\par\vspace{20pt}
\noindent\textcolor{tealshade}{\textbf{Как пользоваться ревью.}} По каждому заданию сначала найдите место расхождения, затем разберите исправление и самостоятельно решите одну новую задачу на закрепление. Ответы к этим новым задачам здесь намеренно не приводятся.

\clearpage
\thispagestyle{empty}
\vspace*{\fill}
\begin{center}
{\Huge\color{tealshade} Разбор ошибок}\par
\vspace{10mm}
\textcolor{tealshade}{\rule{0.56\linewidth}{0.5pt}}
\end{center}
\vspace*{\fill}

\taskheading{9}{Меньшая диагональ ромба}
\subfield{Текст задания.}
Острый угол ромба равен $70^\circ$. Найдите угол между его стороной и меньшей диагональю.

\subfield{Ваше решение.}
В работе записано $70^\circ:2=35^\circ$.

\subfield{Ваш ответ.}
$35^\circ$.

\errorfield
Вы использовали верное свойство диагонали, но выбрали не тот угол ромба.

\subfield{Где именно возникает ошибка.}
Последняя верная идея — диагональ ромба делит угол при своей вершине пополам. Первый неверный шаг — деление данного острого угла $70^\circ$ на два, хотя в условии требуется угол со \emph{меньшей} диагональю. Меньшая диагональ выходит из вершины \emph{тупого}, а не острого угла ромба.

\subfield{Почему это неверно.}
У ромба две диагонали. Обе делят соответствующие углы пополам, но половины этих углов различны. Большая диагональ соединяет вершины острых углов, и угол между ней и стороной действительно составляет $70^\circ:2=35^\circ$. Меньшая диагональ соединяет вершины тупых углов. Поэтому полученное число относится к другой диагонали и не отвечает на вопрос.

\ideafield
Сначала определите, из каких вершин проведена нужная диагональ, и только затем используйте свойство биссектрисы. В ромбе меньшая диагональ соединяет вершины тупых углов.

\subfield{Исправление от последнего правильного шага.}
Свойство «диагональ делит угол пополам» оставляем. Исправляем выбранный угол: соседний с острым углом ромба является тупым и равен
\[
180^\circ-70^\circ=110^\circ.
\]
Теперь делим пополам именно его:
\[
110^\circ:2=55^\circ.
\]

\subfield{Полное правильное решение.}
Обозначим ромб $ABCD$ так, чтобы $\angle A=70^\circ$. Соседние углы параллелограмма, а значит и ромба, в сумме дают $180^\circ$. Следовательно,
\[
\angle B=180^\circ-70^\circ=110^\circ.
\]
Диагональ $BD$ соединяет вершины тупых углов и является меньшей диагональю ромба. Поскольку диагонали ромба делят углы пополам, $BD$ делит $\angle ABC$ на два равных угла:
\[
\angle ABD=\angle DBC=\frac{110^\circ}{2}=\boxed{55^\circ}.
\]
Искомый угол между стороной $BA$ и меньшей диагональю $BD$ равен $55^\circ$.

\subfield{Почему диагональ $BD$ именно меньшая.}
Сравним треугольники $ABD$ и $ABC$. В каждом две стороны равны стороне ромба. При этом сторона $BD$ лежит напротив угла $70^\circ$, а сторона $AC$ — напротив угла $110^\circ$. При одинаковых двух сторонах большему углу между ними соответствует большая третья сторона. Поэтому $AC>BD$.

\subfield{Проверка.}
Половина тупого угла — $55^\circ$, половина острого — $35^\circ$. Эти два угла в сумме равны $90^\circ$, что согласуется с перпендикулярностью диагоналей ромба.

\resultfield
Верный ответ — $\boxed{55^\circ}$. В задачах на диагонали важно сначала различать большую и меньшую диагональ.

\tealheading{Задача на закрепление}
Острый угол ромба равен $64^\circ$. Найдите угол между его стороной и меньшей диагональю.

\taskheading{10}{Разность квадратов и знаки скобок}
\subfield{Текст задания.}
Решите уравнение $x^8=(12-x)^4$ разложением на множители.

\subfield{Ваше решение.}
В работе верно сделан переход
\[
x^8-(12-x)^4=0,
\]
а затем применена разность квадратов:
\[
\bigl(x^4-(12-x)^2\bigr)
\bigl(x^4+(12-x)^2\bigr)=0.
\]
Дальнейший ход содержит попытку разложить первый множитель ещё раз. В заключительной строке получено выражение с множителем $x^2-12-x$; записаны значения $x_1=4$ и $x_2=-3$. Следующая запись корня не завершена.

\subfield{Ваш ответ.}
Записаны числа $4$ и $-3$; итоговая запись не доведена до конца.

\errorfield
При снятии скобок после знака минус изменён не тот знак. Кроме того, не завершён анализ второго множителя.

\subfield{Где именно возникает ошибка.}
До выражения
\[
\bigl(x^2-(12-x)\bigr)\bigl(x^2+(12-x)\bigr)
\]
преобразования первого большого множителя верны. Ошибка появляется в переходе
\[
\underbrace{x^2-(12-x)}_{\text{верно}}
\quad\longrightarrow\quad
\underbrace{x^2-12-x}_{\text{неверно}}.
\]
Знак перед $x$ должен стать плюсом, а не остаться минусом.

\subfield{Почему это неверно.}
Минус перед скобкой означает умножение \emph{каждого} слагаемого внутри неё на $-1$:
\[
-(12-x)=-12+x.
\]
Поэтому
\[
x^2-(12-x)=x^2-12+x=x^2+x-12.
\]
Если написать $x^2-12-x$, получится уже другое квадратное уравнение, у которого действительно корни $4$ и $-3$. Но эти числа не обязаны удовлетворять исходному уравнению — и не удовлетворяют.

\ideafield
Применяйте формулу $a^2-b^2=(a-b)(a+b)$ поэтапно. При раскрытии $-(a-b)$ оба знака внутри скобки меняются. После разложения произведения на множители проверяйте \emph{каждый} множитель на нуль, а также исключайте те, которые никогда не равны нулю.

\subfield{Исправление от последнего правильного шага.}
Вернёмся к правильной записи:
\[
\bigl(x^4-(12-x)^2\bigr)
\bigl(x^4+(12-x)^2\bigr)=0.
\]
Разложим только первый множитель:
\begin{align*}
x^4-(12-x)^2
&=\bigl(x^2-(12-x)\bigr)\bigl(x^2+(12-x)\bigr)\\
&=(x^2+x-12)(x^2-x+12).
\end{align*}
Второй большой множитель $x^4+(12-x)^2$ тоже учитываем: он положителен при любом действительном $x$, а значит нулём быть не может.

\subfield{Полное правильное решение.}
Перенесём всё в левую часть и применим формулу разности квадратов:
\begin{align*}
x^8-(12-x)^4&=0,\\
\bigl(x^4-(12-x)^2\bigr)
\bigl(x^4+(12-x)^2\bigr)&=0.
\end{align*}

\textbf{Первый случай:} $x^4+(12-x)^2=0$.
Оба слагаемых — квадраты и неотрицательны. Их сумма могла бы равняться нулю лишь при одновременном выполнении $x=0$ и $12-x=0$. Это невозможно. Следовательно, решений в этом случае нет.

\textbf{Второй случай:} $x^4-(12-x)^2=0$.
Ещё раз применяем разность квадратов и раскрываем скобки:
\begin{align*}
\bigl(x^2-(12-x)\bigr)\bigl(x^2+(12-x)\bigr)&=0,\\
(x^2+x-12)(x^2-x+12)&=0.
\end{align*}
Теперь рассмотрим два квадратных уравнения.

\textbf{Первое квадратное уравнение:}
\[
x^2+x-12=0.
\]
Раскладываем трёхчлен на множители:
\[
(x+4)(x-3)=0,
\qquad x=-4\ \text{или}\ x=3.
\]

\textbf{Второе квадратное уравнение:}
\[
x^2-x+12=0.
\]
Его дискриминант
\[
D=(-1)^2-4\cdot1\cdot12=1-48=-47<0.
\]
Действительных корней нет. Значит, все действительные решения исходного уравнения: $-4$ и $3$.

\subfield{Проверка.}
При $x=3$ получаем
\[
3^8=(12-3)^4=9^4=(3^2)^4=3^8.
\]
При $x=-4$ получаем
\[
(-4)^8=(12-(-4))^4=16^4=(4^2)^4=4^8.
\]
Оба значения удовлетворяют исходному уравнению. Числа из неверного решения не проходят проверку: при $x=4$ имеем $4^8\ne8^4$, а при $x=-3$ — $(-3)^8\ne15^4$.

\resultfield
Ответ: $\boxed{x=-4\ \text{или}\ x=3}$. Важно не только правильно применять разность квадратов, но и аккуратно раскрывать скобки и исследовать все полученные множители.

\tealheading{Задача на закрепление}
Решите уравнение $x^8=(20-x)^4$ разложением на множители.

\vfill
\begin{center}
\textcolor{tealshade}{\rule{0.42\linewidth}{0.45pt}}\\[6pt]
{\footnotesize Лёвин Артём Александрович эксклюзивно для Ярослава Верещагина}
\end{center}
\end{document}